\documentclass[tikz,border=0pt]{standalone}
\usepackage[UTF8]{ctex}
\usepackage{xcolor}
\usetikzlibrary{arrows.meta,calc,positioning,shapes.geometric}

\definecolor{accent}{HTML}{FF3CAC}

\tikzset{
  box/.style={
    draw=white,
    text=white,
    line width=0.8pt,
    align=center,
    font=\sffamily\fontsize{7.5}{9.5}\selectfont,
    inner sep=5pt
  },
  decision/.style={
    diamond,
    aspect=3.7,
    draw=white,
    text=white,
    line width=0.8pt,
    align=center,
    font=\sffamily\fontsize{7.5}{9.5}\selectfont,
    inner sep=2pt
  },
  flow/.style={
    draw=white,
    line width=0.8pt,
    -{Stealth[length=2.8mm,width=2mm]}
  },
  retry/.style={
    draw=white,
    line width=0.8pt,
    dashed,
    -{Stealth[length=2.8mm,width=2mm]}
  },
  zone/.style={
    draw=white,
    line width=0.6pt,
    dash pattern=on 5pt off 4pt
  },
  accent text/.style={
    text=accent,
    font=\sffamily\bfseries\fontsize{11}{14}\selectfont
  },
  note text/.style={
    text=accent,
    font=\sffamily\bfseries\fontsize{7.5}{9.5}\selectfont
  }
}

\begin{document}
\begin{tikzpicture}[x=1cm,y=1cm]
  \path[use as bounding box] (0,0) rectangle (15.36,10.24);
  \fill[black] (0,0) rectangle (15.36,10.24);

  \draw[zone] (0.20,8.65) -- (15.16,8.65);
  \draw[zone] (0.20,2.50) -- (15.16,2.50);
  \node[accent text,anchor=east] at (15.05,8.98) {用户态};
  \node[accent text,anchor=east] at (15.05,8.32) {内核态};
  \node[accent text,anchor=east] at (15.05,2.18) {硬件};

  \node[box,text width=4.10cm,minimum height=0.85cm] (process) at (2.90,9.55)
    {{\color{accent}\bfseries 进程 P}\\[2pt] \texttt{write(1, buffer, count)}\\[-1pt]
     \fontsize{6.5}{8}\selectfont （标准输出连接终端）};

  \node[box,text width=4.10cm,minimum height=0.68cm] (lookup) at (2.90,8.00)
    {由 \texttt{fd=1} 找到终端设备\\对应的写操作};

  \draw[box] (0.65,6.80) -- (2.90,7.35) -- (5.15,6.80) -- (2.90,6.25) -- cycle;
  \node[text=white,align=center,font=\sffamily\fontsize{7.5}{9.5}\selectfont] at (2.90,6.80)
    {输出缓冲区\\有足够空间？};
  \coordinate (space-north) at (2.90,7.35);
  \coordinate (space-south) at (2.90,6.25);
  \coordinate (space-east) at (5.15,6.80);

  \node[box,text width=4.10cm,minimum height=0.68cm] (copy) at (2.90,5.42)
    {从用户缓冲区复制数据};

  \node[box,text width=4.10cm,minimum height=0.68cm] (kbuf) at (2.90,4.32)
    {内核输出缓冲区};

  \node[box,text width=2.55cm,minimum height=0.68cm] (return) at (1.78,3.18)
    {系统调用返回};
  \node[note text,anchor=west] at (0.32,2.68) {不等待物理输出完成};

  \node[box,text width=3.55cm,minimum height=1.05cm] (block) at (7.55,6.80)
    {缓冲区空间不足\\阻塞 P，加入输出等待队列\\调度其他进程};

  \node[box,text width=3.15cm,minimum height=0.68cm] (ready) at (7.55,5.18)
    {P 进入就绪队列};

  \node[box,text width=3.15cm,minimum height=0.68cm] (rerun) at (7.55,3.72)
    {调度器以后再次运行 P};

  \node[box,text width=3.70cm,minimum height=1.18cm] (irq) at (12.42,5.18)
    {{\color{accent}\bfseries 设备输出中断处理程序}\\[3pt]
     \fontsize{7}{9}\selectfont 确认发送完成，继续发送\\
     缓冲区有空间时唤醒 P};

  \node[box,text width=3.70cm,minimum height=0.90cm] (driver) at (12.42,3.45)
    {{\color{accent}\bfseries 终端驱动程序}\\[3pt] 把缓冲区数据交给控制器};

  \node[box,text width=3.05cm,minimum height=0.72cm] (controller) at (12.42,1.68)
    {终端控制器发送数据\\[-1pt]\fontsize{7}{9}\selectfont 就绪后产生中断};

  \node[box,text width=3.05cm,minimum height=0.62cm] (device) at (12.42,0.58)
    {显示器或串口};

  \draw[flow] (process.south) -- (lookup.north);
  \draw[flow] (lookup.south) -- (space-north);
  \draw[flow] (space-south) -- node[right,note text] {YES} (copy.north);
  \draw[flow] (copy.south) -- (kbuf.north);

  \draw[flow] (space-east) -- node[above,note text] {NO} (block.west);

  \draw[flow] (kbuf.south) -- ++(0,-0.22) -| (return.north);
  \draw[flow] (return.west) -- (0.22,3.18) -- (0.22,9.55) -- (process.west);

  \draw[flow] (kbuf.east) -- (5.55,4.32) -- (5.55,2.82) --
    (10.15,2.82) -- (10.15,3.45) -- (driver.west);
  \draw[flow] (driver.south) -- (controller.north);
  \draw[flow] (controller.south) -- (device.north);

  \draw[flow] (controller.east) -- (14.85,1.68) --
    node[right,note text,rotate=90] {设备中断} (14.85,5.18) -- (irq.east);
  \draw[flow] (irq.west) -- (ready.east);
  \draw[flow] (ready.south) -- (rerun.north);

  \draw[retry] (rerun.west) -- (5.42,3.72) -- (5.42,6.80) -- (space-east);
\end{tikzpicture}
\end{document}
